\documentclass[11pt,a4paper]{article}

% Packages
\usepackage[utf8]{inputenc}
\usepackage[T1]{fontenc}
\usepackage{geometry}
\usepackage{listings}
\usepackage{xcolor}
\usepackage{booktabs}
\usepackage{hyperref}
\usepackage{fancyhdr}
\usepackage{graphicx}
\usepackage{tcolorbox}
\usepackage{enumitem}

% Page geometry
\geometry{margin=1in}

% Colors
\definecolor{codegreen}{rgb}{0,0.6,0}
\definecolor{codegray}{rgb}{0.5,0.5,0.5}
\definecolor{codepurple}{rgb}{0.58,0,0.82}
\definecolor{backcolour}{rgb}{0.95,0.95,0.92}
\definecolor{tipblue}{rgb}{0.9,0.95,1}
\definecolor{warnred}{rgb}{1,0.95,0.9}

% Code listings style
\lstdefinestyle{bashstyle}{
    backgroundcolor=\color{backcolour},
    commentstyle=\color{codegreen},
    keywordstyle=\color{codepurple},
    numberstyle=\tiny\color{codegray},
    stringstyle=\color{codepurple},
    basicstyle=\ttfamily\footnotesize,
    breakatwhitespace=false,
    breaklines=true,
    captionpos=b,
    keepspaces=true,
    numbers=none,
    showspaces=false,
    showstringspaces=false,
    showtabs=false,
    tabsize=2,
    frame=single,
    framerule=0.5pt,
}

\lstdefinestyle{yamlstyle}{
    backgroundcolor=\color{backcolour},
    basicstyle=\ttfamily\footnotesize,
    breaklines=true,
    frame=single,
    framerule=0.5pt,
}

\lstset{style=bashstyle}

% Header/Footer
\pagestyle{fancy}
\fancyhf{}
\rhead{ModelSpec Introspector}
\lhead{GCP Testing Guide}
\rfoot{Page \thepage}

% Hyperref setup
\hypersetup{
    colorlinks=true,
    linkcolor=blue,
    urlcolor=blue,
}

\begin{document}

% Title
\begin{center}
    {\LARGE\textbf{ModelSpec Introspector}}\\[0.5em]
    {\Large GCP Testing Guide}\\[1em]
    \textit{Estimated Total Cost: \$2--5 for complete feature testing}
\end{center}

\tableofcontents
\newpage

%==============================================================================
\section{Recommended Model \& GPU Selection}
%==============================================================================

\subsection{Best Choice: TinyLlama-1.1B on T4 GPU}

\begin{table}[h]
\centering
\begin{tabular}{@{}lll@{}}
\toprule
\textbf{Aspect} & \textbf{Recommendation} & \textbf{Why} \\
\midrule
Model & \texttt{TinyLlama/TinyLlama-1.1B-Chat-v1.0} & Only 2GB VRAM, fast inference \\
GPU & NVIDIA T4 (16GB) & Cheapest GPU (\textasciitilde\$0.35/hr Spot) \\
Machine & \texttt{n1-standard-4} & 4 vCPU, 15GB RAM \\
Region & \texttt{us-central1-a} & Best availability \\
\bottomrule
\end{tabular}
\caption{Recommended configuration for cost-effective testing}
\end{table}

\subsection{Alternative Models}

\begin{table}[h]
\centering
\begin{tabular}{@{}llll@{}}
\toprule
\textbf{Model} & \textbf{Parameters} & \textbf{VRAM Needed} & \textbf{GPU Required} \\
\midrule
TinyLlama-1.1B & 1.1B & \textasciitilde2GB & T4 \checkmark \\
Phi-2 & 2.7B & \textasciitilde6GB & T4 \checkmark \\
Mistral-7B (INT4) & 7B & \textasciitilde8GB & T4 \checkmark \\
Mistral-7B (FP16) & 7B & \textasciitilde16GB & L4 or A100 \\
\bottomrule
\end{tabular}
\caption{Alternative model options}
\end{table}

%==============================================================================
\section{GCP Setup (Step-by-Step)}
%==============================================================================

\subsection{Step 1: Enable APIs \& Set Project}

\begin{lstlisting}[language=bash]
# Set your project
gcloud config set project YOUR_PROJECT_ID

# Enable required APIs
gcloud services enable container.googleapis.com
gcloud services enable compute.googleapis.com
\end{lstlisting}

\subsection{Step 2: Create GKE Cluster with T4 GPU}

\begin{lstlisting}[language=bash]
# Create a minimal GKE cluster with Spot instances
gcloud container clusters create modelspec-test \
  --zone=us-central1-a \
  --machine-type=n1-standard-4 \
  --accelerator=type=nvidia-tesla-t4,count=1 \
  --num-nodes=1 \
  --spot \
  --disk-size=50GB \
  --enable-autoscaling \
  --min-nodes=0 \
  --max-nodes=1

# Install NVIDIA GPU drivers
kubectl apply -f https://raw.githubusercontent.com/GoogleCloudPlatform/\
container-engine-accelerators/master/nvidia-driver-installer/\
cos/daemonset-preloaded.yaml

# Get credentials
gcloud container clusters get-credentials modelspec-test \
  --zone=us-central1-a
\end{lstlisting}

\subsection{Step 3: Verify GPU is Available}

\begin{lstlisting}[language=bash]
# Wait ~2 minutes for driver installation, then:
kubectl get nodes -o jsonpath='{.items[*].status.allocatable.nvidia\.com/gpu}'
# Should output: 1
\end{lstlisting}

%==============================================================================
\section{Deploy vLLM with TinyLlama}
%==============================================================================

\subsection{Create Namespace}

\begin{lstlisting}[language=bash]
kubectl create namespace inference
\end{lstlisting}

\subsection{Create vLLM Deployment}

Save the following as \texttt{vllm-tinyllama.yaml}:

\begin{lstlisting}[style=yamlstyle]
apiVersion: apps/v1
kind: Deployment
metadata:
  name: vllm-tinyllama
  namespace: inference
  labels:
    app: vllm-tinyllama
    app.kubernetes.io/name: vllm
    framework: vllm
spec:
  replicas: 1
  selector:
    matchLabels:
      app: vllm-tinyllama
  template:
    metadata:
      labels:
        app: vllm-tinyllama
        app.kubernetes.io/name: vllm
        framework: vllm
    spec:
      containers:
      - name: vllm
        image: vllm/vllm-openai:latest
        args:
          - "--model"
          - "TinyLlama/TinyLlama-1.1B-Chat-v1.0"
          - "--served-model-name"
          - "tinyllama"
          - "--dtype"
          - "float16"
          - "--max-model-len"
          - "2048"
          - "--tensor-parallel-size"
          - "1"
          - "--gpu-memory-utilization"
          - "0.9"
          - "--enable-prefix-caching"
        ports:
          - containerPort: 8000
            name: http
        resources:
          requests:
            nvidia.com/gpu: 1
            cpu: "2"
            memory: "8Gi"
          limits:
            nvidia.com/gpu: 1
            cpu: "4"
            memory: "12Gi"
        readinessProbe:
          httpGet:
            path: /health
            port: 8000
          initialDelaySeconds: 120
          periodSeconds: 10
        livenessProbe:
          httpGet:
            path: /health
            port: 8000
          initialDelaySeconds: 180
          periodSeconds: 30
---
apiVersion: v1
kind: Service
metadata:
  name: vllm-tinyllama
  namespace: inference
spec:
  selector:
    app: vllm-tinyllama
  ports:
    - port: 8000
      targetPort: 8000
      name: http
  type: ClusterIP
\end{lstlisting}

\subsection{Deploy and Verify}

\begin{lstlisting}[language=bash]
# Deploy
kubectl apply -f vllm-tinyllama.yaml

# Watch pod status (wait for Running)
kubectl get pods -n inference -w

# Check logs for model loading progress
kubectl logs -n inference -l app=vllm-tinyllama -f

# Port-forward to test
kubectl port-forward -n inference svc/vllm-tinyllama 8000:8000 &

# Test the API
curl http://localhost:8000/v1/models
curl http://localhost:8000/health
curl http://localhost:8000/metrics | head -50
\end{lstlisting}

%==============================================================================
\section{Install \& Test ModelSpec Introspector}
%==============================================================================

\subsection{Install the Tool}

\begin{lstlisting}[language=bash]
# Clone and install
cd /path/to/ModelSpec
poetry install

# Or install directly
pip install -e .
\end{lstlisting}

\subsection{Apply RBAC Permissions}

\begin{lstlisting}[language=bash]
kubectl apply -f rbac/
\end{lstlisting}

\subsection{Run the Tests}

\subsubsection{Test 1: Connection Test}
\begin{lstlisting}[language=bash]
piqc test-connection
\end{lstlisting}
\textbf{Expected:} Connected to cluster

\subsubsection{Test 2: Basic Scan (Table Output)}
\begin{lstlisting}[language=bash]
piqc scan --namespace inference
\end{lstlisting}
\textbf{Expected:} Table showing vllm-tinyllama deployment

\subsubsection{Test 3: YAML Output}
\begin{lstlisting}[language=bash]
piqc scan --namespace inference --output yaml \
  > output/test-modelspec.yaml
\end{lstlisting}
\textbf{Expected:} YAML file with complete ModelSpec

\subsubsection{Test 4: JSON Output}
\begin{lstlisting}[language=bash]
piqc scan --namespace inference --output json \
  > output/test-modelspec.json
\end{lstlisting}

\subsubsection{Test 5: PIQC Facts Bundle}
\begin{lstlisting}[language=bash]
piqc scan --namespace inference --output-piqc \
  > output/piqc-facts.json
\end{lstlisting}
\textbf{Expected:} JSON with facts like \texttt{runtime.engineType}, \texttt{vllm.dtype}

\subsubsection{Test 6: GPU Metrics Collection}
\begin{lstlisting}[language=bash]
piqc scan --namespace inference --output yaml -v
\end{lstlisting}
\textbf{Expected:} GPU info with utilization, memory, temperature

\subsubsection{Test 7: Runtime Metrics (Prometheus)}
\begin{lstlisting}[language=bash]
piqc scan --namespace inference --collect-runtime \
  --output yaml
\end{lstlisting}
\textbf{Expected:} \texttt{runtimeState} section with live metrics

\subsubsection{Test 8: Multi-Namespace Scan}
\begin{lstlisting}[language=bash]
piqc scan --output yaml
\end{lstlisting}
\textbf{Expected:} Scans all namespaces, finds inference workloads

\subsubsection{Test 9: Combined Output}
\begin{lstlisting}[language=bash]
piqc scan --combined --output yaml \
  > output/all-modelspecs.yaml
\end{lstlisting}

%==============================================================================
\section{Feature Verification Checklist}
%==============================================================================

\begin{table}[h]
\centering
\small
\begin{tabular}{@{}lll@{}}
\toprule
\textbf{Feature} & \textbf{How to Verify} & \textbf{CLI Flag} \\
\midrule
vLLM Detection & Check \texttt{engine.name: vllm} & Default \\
Model Inference & Check \texttt{model.architecture: llama} & Default \\
GPU Metrics & Check \texttt{resources.gpus[].utilization} & Default \\
Precision Detection & Check \texttt{inference.precision: float16} & Default \\
Max Model Len & Check \texttt{inference.maxModelLen: 2048} & Default \\
Tensor Parallel & Check \texttt{inference.tensorParallelSize: 1} & Default \\
Runtime Metrics & Check \texttt{runtimeState.vllm.gpuCacheUsagePercent} & \texttt{--collect-runtime} \\
PIQC Facts & Check \texttt{runtime.engineType}, \texttt{hardware.gpuType} & \texttt{--output-piqc} \\
Pod Labels & Check \texttt{metadata.labels} captured & Default \\
Namespace Filter & Only scan specific namespace & \texttt{--namespace} \\
\bottomrule
\end{tabular}
\caption{Feature verification checklist}
\end{table}

%==============================================================================
\section{Cleanup (IMPORTANT -- Stop Charges!)}
%==============================================================================

\begin{tcolorbox}[colback=warnred,colframe=red!50!black,title=\textbf{Important: Cleanup to Stop Charges}]
Always delete your cluster when testing is complete to avoid ongoing charges.
\end{tcolorbox}

\begin{lstlisting}[language=bash]
# Delete the deployment first
kubectl delete -f vllm-tinyllama.yaml

# Delete the cluster (stops all charges)
gcloud container clusters delete modelspec-test \
  --zone=us-central1-a --quiet

# Verify deletion
gcloud container clusters list
\end{lstlisting}

%==============================================================================
\section{Cost Breakdown}
%==============================================================================

\begin{table}[h]
\centering
\begin{tabular}{@{}lll@{}}
\toprule
\textbf{Resource} & \textbf{Hourly Cost} & \textbf{For 2hr Test} \\
\midrule
T4 GPU (Spot) & \textasciitilde\$0.35 & \$0.70 \\
n1-standard-4 (Spot) & \textasciitilde\$0.05 & \$0.10 \\
GKE Management Fee & \textasciitilde\$0.10 & \$0.20 \\
\midrule
\textbf{Total} & & \textbf{\textasciitilde\$1--2} \\
\bottomrule
\end{tabular}
\caption{Estimated testing costs}
\end{table}

\begin{tcolorbox}[colback=tipblue,colframe=blue!50!black,title=\textbf{Pro Tip}]
Use the \texttt{--spot} flag when creating your cluster to save 60--70\% on GPU costs!
\end{tcolorbox}

%==============================================================================
\section{Troubleshooting}
%==============================================================================

\subsection{Pod Stuck in Pending}
\begin{lstlisting}[language=bash]
kubectl describe pod -n inference -l app=vllm-tinyllama
\end{lstlisting}
Usually means GPU quota not available. Try a different zone.

\subsection{GPU Driver Not Installing}
\begin{lstlisting}[language=bash]
kubectl logs -n kube-system -l k8s-app=nvidia-driver-installer
\end{lstlisting}

\subsection{Model Loading Timeout}
The readiness probe has a 2-minute timeout. TinyLlama loads in approximately 1 minute.

\subsection{Permission Denied Errors}
\begin{lstlisting}[language=bash]
kubectl apply -f rbac/
\end{lstlisting}
Ensure RBAC is applied before scanning.

%==============================================================================
\section{Quick Start (Copy-Paste Version)}
%==============================================================================

\begin{lstlisting}[language=bash]
# 1. Create cluster
gcloud container clusters create modelspec-test \
  --zone=us-central1-a \
  --machine-type=n1-standard-4 \
  --accelerator=type=nvidia-tesla-t4,count=1 \
  --num-nodes=1 --spot

# 2. Install GPU drivers
kubectl apply -f https://raw.githubusercontent.com/GoogleCloudPlatform/\
container-engine-accelerators/master/nvidia-driver-installer/\
cos/daemonset-preloaded.yaml

# 3. Get credentials
gcloud container clusters get-credentials modelspec-test \
  --zone=us-central1-a

# 4. Create namespace and deploy
kubectl create namespace inference
kubectl apply -f vllm-tinyllama.yaml

# 5. Wait for pod (2-3 minutes)
kubectl wait --for=condition=ready pod -l app=vllm-tinyllama \
  -n inference --timeout=300s

# 6. Apply RBAC and scan
kubectl apply -f rbac/
piqc scan --namespace inference --output yaml

# 7. CLEANUP (don't forget!)
gcloud container clusters delete modelspec-test \
  --zone=us-central1-a --quiet
\end{lstlisting}

\end{document}
